\section{Research status, corrections, and further questions}
\label{mld:sec:research}

\subsection{What the continuation settles}

The main resolved question in the supplied incoming set concerns integer
dilation and the choice of local subtraction coordinate. Theorem~
\ref{mld:dil:scale-thm} gives the complete local correction, including every
Stieltjes index and argument derivative. Equations~\eqref{mld:dil:full-formula}
and \eqref{mld:dil:all-derivative-formula} then evaluate unequal-dilation
products when the singular grids are disjoint. Their dependence on
$\log\lcm(p,q)$ is explicit. These statements answer the integer-covering
case of the rescaling question; a general nonlinear coordinate change
is outside their scope.

The exact locations are the further-research items on unequal dilations
in \emph{Bilinear Closure of Generalized Stieltjes Correlations}
\cite{IncomingClosure}, on dilation, traces and subtraction scale in
\emph{Shifted Hurwitz-Zeta Correlations} \cite{IncomingShifted}, and on
rescaling in \emph{A Convolution Calculus for Stieltjes Functions}
\cite{IncomingCalculus}. The integration notes give the archive members
and line locations in the pinned snapshot.

The Mellin results are independent of that regularization problem.
They give ordinary integrals with a uniform finite formula at arbitrary
denominator order. The integer resonances, all their logarithmic moments,
and the spectral derivatives in Section~\ref{mld:sec:jets} are proved
identities. The harmonic reflection is likewise an identity of
holomorphic functions and ordinary convergent series. Its even-weight
unshifted coefficient evaluations have published antecedents
\cite{XuWang}; the two-shift formulation is presented with a proof and
without an unverified claim of historical priority.

\begin{center}
\begin{tabularx}{\textwidth}{@{}>{\raggedright\arraybackslash}p{0.35\textwidth}Y@{}}
\toprule
Claim & Status after this work\\
\midrule
All-$N$ Mellin transform and integer resonances & Proved in
Theorems~\ref{mld:thm:Mellin-master} and \ref{mld:thm:integer-resonance}.\\
Shifted generalized-harmonic reflection & Proved in
Theorem~\ref{mld:gh:thm:reflection}; scalar parity precedent credited.\\
Integer-dilation finite parts and disjoint-grid products & Proved in
Theorems~\ref{mld:dil:scale-thm}, \ref{mld:dil:full-stieltjes}, and
\ref{mld:dil:all-derivatives}.\\
Undilated Stieltjes closure and convolution & Already proved in the
supplied incoming material; recalled as inputs.\\
Gaussian $S_4$ evaluation & Already proved in the pinned manuscript;
its status is unchanged.\\
Current Gaussian $S_6$ and $S_8$ evaluations & Still conjectural.
No equality follows from a small residual.\\
Colliding unequal singular grids; nonlinear coordinate changes &
Not treated by the new dilation theorem.\\
\bottomrule
\end{tabularx}
\end{center}

\subsection{The Gaussian conjectures and the missing parity component}

The manuscript defines
\[
 S_p=\sum_{n\ge0}\frac{(-1)^nH_n}{(2n+1)^p},\qquad
 g_{a,b}=\operatorname{Im}\Li_{a,b}(\ii,1),
\]
where
$\Li_{a,b}(x,y)=\sum_{n>m\ge1}x^ny^m/(n^am^b)$ in its convergence
region. Its proved translation is
\[
 S_p=\operatorname{Im}\bigl(\Li_{p,1}(\ii,1)
                                  +\Li_{p,1}(\ii,-1)\bigr).
\]
The file \texttt{04-S4-proof.tex} supplies the exact $S_4$ certificate
and a convergent octahedral relation. The current $S_6$ conjecture is in
\texttt{04-cyclotomic-quotients.tex}; the current $S_8$ conjecture is in
\texttt{04-S8-candidate.tex}. Their complete coefficient vectors and
source labels are preserved in \texttt{integration/GAUSSIAN\_STATUS.md}.

For specificity, the $S_6$ equality still to be proved is
\begin{align}
 S_6\stackrel{?}{={}}&\frac{722}{527}g_{6,1}
  +\frac{40}{527}g_{4,3}-\frac{128}{155}g_{2,5}
  +\frac{15191}{28569600}\pi^7\notag\\
 &-\frac3{124}\beta(2)\zeta(5)
  -\frac{2373}{10540}\beta(4)\zeta(3)-2\beta(6)\log2.
 \label{mld:eq:S6-open}
\end{align}
The manuscript also proves a change of coordinates between two
presentations of this candidate. That identity equates two residual
expressions; it does not show that either residual is zero.

The obstruction to proving \eqref{mld:eq:S6-open} by the reflection obtained
here is exact. Its Taylor coefficient is
\[
 [z^{p-1}u^{r-1}]\{A(z,u)+A(-z,-u)\}
           =(1+(-1)^{p+r})Q_{p,r}.
\]
Since $S_6=Q_{6,1}$ and $S_8=Q_{8,1}$ have total weights $7$ and $9$,
their coefficients vanish in this projection. The same is true of the
weight-$5$ $S_4$ value: its known proof uses additional functional
information. Further differentiation of the same even projection
does not recover its absent Taylor coefficients. A successful extension
must control the complementary function
$A(z,u)-A(-z,-u)$ or produce an independent relation among the required
mixed-color periods.

The separating functional already exhibited in the manuscript rules out
a particular restricted row-span certificate for the $S_6$ candidate.
It does not rule out the identity itself. Similarly, a formal quotient
dimension does not prove arithmetic independence of the numerical
periods. These distinctions are retained throughout this continuation.

\subsection{Audit findings and concrete integration corrections}

No new false mathematical statement was found in the sampled pinned
Gaussian claims or the incoming correlation identities used here.
Several concrete reconciliations are nevertheless needed when the
packages are integrated.

\begin{enumerate}[label=\arabic*.]
\item \textbf{Reconcile the open-question lists.}
The Closure report's questions about undilated parameter derivatives and
associative convolution are answered in the current Shifted Hurwitz and
Convolution packages. Their entries should point to those results.
The integer-dilation questions should point to the present scale and
product theorems, with disjointness and ambient normalization stated.

\item \textbf{Keep the two dilation operations distinct.}
The replacement $X_{q,n,r}=U_qT_{n,r}$ would omit
\eqref{mld:dil:scale-law}. A simple exact diagnostic is
\[
 \langle U_qG_0,1\rangle=0,\qquad
 \langle X_{q,0,0},1\rangle=-\log q.
\]
Thus they already differ at zero Stieltjes index and zero derivative
order. This corrects a possible naive extension of the incoming formulas;
it does not assert that the supplied undilated theorem made that error.

\item \textbf{State the cutoff variable in derivative identities.}
The valid undilated assertion $h_{r,n}=0$ for $n\ge r$ cannot be copied
unchanged into a theorem about a fixed ambient scale. For example,
\[
 DX_{q,1,0}=qX_{q,1,1}-\frac{\log q}{q}\Delta_q'.
\]
A separate constant term in each local variable has the canonical
normalization used here. Rescaling those variables before one common
constant-term extraction changes it.

\item \textbf{Retain resonance values before simplifying factors.}
At integer spectral parameters, use the analytic germ
\eqref{mld:eq:cancelled-zeta-germ}. The equality $F_0(1)=1$ is responsible
for the $-1/4$ and $-3/4$ higher-pole examples. Dropping a factor that
vanishes next to a zeta pole is not a valid simplification.

\item \textbf{Preserve the existing conjecture and error history.}
The present $S_8$ candidate is different from the earlier vector already
rejected in \texttt{10-discovery.tex}. The mixed-color symmetry
counterexample is already recorded in \texttt{04-depth.tex}.
These are existing corrections, not new discoveries of this audit.
The generalized-harmonic family $Q_{p,r}$ should also retain its own
notation, since the manuscript's $S_{p,r}$ uses elementary symmetric
harmonic sums.
\end{enumerate}

\subsection{Proposed research programme}

The following questions seek exact identities or structural reductions.
They are research questions, not conjectures inferred from a few
numerical samples. The first three are the most immediate next steps.

\begin{enumerate}[label=\textbf{Q\arabic*.},leftmargin=2.7em]
\item \textbf{The missing Gaussian projection.}
Find a functional equation or convergent iterated-integral relation that
determines the weight-$7$ and weight-$9$ coefficients of
$A(z,u)-A(-z,-u)$. A concrete success criterion is an exact zero
certificate for \eqref{mld:eq:S6-open}, followed by the current $S_8$ vector.
The known $S_4$ octahedral seed suggests testing genuinely new functional
relations before enlarging a numerically fitted coefficient basket.

\item \textbf{Colliding unequal dilations.}
Extend the ambient-coordinate product theorem to $Pa\in\Z$. There the
partition separating the two singular grids fails, and the product
contains higher singular powers at the common points. Derive the full
local subtraction polynomial and compare its finite value with the
collision expansion of \eqref{mld:dil:full-formula}. The already known
undilated coincident closure is a boundary case and a normalization
check, not a solution of every unequal-grid configuration.

\item \textbf{Nonlinear changes of local coordinate.}
Replace $qt$ near a singularity by an analytic map
$f(t)=qt+c_2t^2+\cdots$, $q>0$. Determine the finite combination of
delta derivatives that converts pullback into ambient finite part for
every $\gamma_n^{(r)}$. The linear coefficient must recover
$c_{n,r}(\log q)$; higher coefficients should be computable from finite
Taylor and Bell-polynomial algebra. The goal is a proved composition
law for the full local correction.

\item \textbf{Three translated or dilated factors.}
For three pairwise disjoint singular grids, determine the spectral
generating kernel of the product of three Hurwitz or Stieltjes
functions. Fourier orthogonality leaves a two-dimensional frequency
constraint. Identify which double Hurwitz sums or multiple
polylogarithms survive, and derive the first nontrivial exact Stieltjes
coefficient formulas. Do not assume the bilinear single-Stieltjes
closure persists at degree three.

\item \textbf{Complete primitive calculus for dilated products.}
Insert the finite closure into the anchored Hurwitz-jet primitive
ladder to construct explicit primitives in the shift $a$ on every
noncollision interval. Then determine the matching constants when
crossing a collision point in a specified finite-part convention.
This would connect local exact integration with a global periodic
distribution statement.

\item \textbf{Character-weighted dilation and multiplication.}
Replace the unweighted trace by Dirichlet-character weights on the
finite covering. Determine how the local correction interacts with
the character's zero mode, and express the resulting correlations in
Dirichlet $L$-jets. The residue-class trace polynomial
\eqref{mld:dil:trace-polynomial} supplies the unweighted test case.

\item \textbf{Rational-shift harmonic identities beyond a reflection.}
Combine the two-shift theorem with recurrences in its inner and outer
shifts. Determine which finite combinations of individual shifted
harmonic sums admit depth-one cyclotomic expressions, rather than only
their reflected pairs. A useful first classification fixes small
denominators in $u$ and $z$ and records exact functional certificates.

\item \textbf{Mellin kernels with noninteger denominator order.}
Replace $N$ in $(1+x)^{-N}$ by a complex parameter. Find a spectral
representation that specializes to the finite polynomial operator at
positive integers, and determine exactly when a finite ordinary-
polylogarithm closure survives. Hypergeometric or generalized Lerch
functions are natural intermediate objects, but their branch and
resonance conventions must be derived from the convergent integral.

\item \textbf{Complex poles and contour rotation.}
Extend Corollary~\ref{mld:cor:rational-kernel} to rational kernels with
nonreal poles. Track the residue terms when a rotated integration ray
crosses a polylogarithm cut. The target is an exact continuation formula
that reduces to the positive-real scaling law when every pole lies on
the negative real axis.

\item \textbf{Products of polylogarithms under the Mellin kernel.}
Study
\[
 \int_0^\infty \frac{x^{a-1}\Li_s(-zx)\Li_t(-wx)}{(1+x)^N}\dd x
\]
where ordinary convergence holds. Decide whether integer Mellin
parameters lead to a finite multiple-polylogarithm algebra, and identify
the minimal depth forced by two independent arguments. The one-factor
integer table gives exact boundary cases and a recurrence in $N$.

\item \textbf{Compatibility of the two harmonic coefficient calculi.}
Both cancelled zeta poles and dilation corrections involve
$\prod_j(1-u/j)$, whereas repeated anchored primitives involve its
reciprocal. Formulate and prove a finite algebra connecting these
operations, including derivative--primitive commutators and their
integration polynomials. This could reduce duplicated coefficient
calculations across the manuscript.

\item \textbf{Exact certificates and formal verification.}
Encode the polynomial recurrences, the scale composition law, and the
integer moment table as exact certificates with stable normalization
metadata. A proof-assistant development could then isolate the finite
algebra from the analytic tasks of dominated parameter differentiation,
meromorphic continuation, and distributional pullback. Numerical
agreement alone should not change theorem status.
\end{enumerate}

These questions are constrained by the proved formulas: each has a
specified base case, normalization, or missing component to resolve.
They also preserve the main distinction of this article: ordinary
polylogarithmic integrals, finite harmonic reductions, and normalized
distribution products are connected by exact operations, but are not
interchangeable without their hypotheses.
